\section{Introduction}
Derived categories play a central role in homological algebra and representation theory, where they provide a framework for comparing algebraic objects beyond their abelian or module categories. In the representation theory of quivers and posets, one is naturally led to derived categories of diagram categories. If $P$ is a finite poset and $\mathcal{A}$ is an abelian category, the category $\mathcal{A}^P$ of $P$-shaped diagrams in $\mathcal{A}$ is again abelian and one may study its derived category $D(\mathcal{A}^P)$. 

For a fixed coefficient category -- for instance, finite-dimensional vector spaces over a field $k$ -- one may ask when two posets give rise to equivalent derived categories of representations. Such an equivalence is, a priori, tied to the chosen coefficients. A stronger phenomenon occurs when the equivalence is completely independent of this choice: two finite posets $P$ and $Q$ are called universally derived equivalent if
\[D(\mathcal{A}^P) \simeq D(\mathcal{A}^Q)\] for every abelian category $\mathcal{A}$. 
Universal derived equivalence is therefore a genuinely combinatorial property of the indexing posets. The problem is then to understand which transformations of a finite poset preserve its derived category in this universal sense. 

One of the main constructions of universally derived equivalent posets is due to Ladkani \cite{ladkani}. Starting from two finite posets together with suitable combinatorial data relating them, his construction produces two different partial orders on their disjoint union and proves that the corresponding diagram categories are universally derived equivalent (see Section \ref{Ladkani's construction}).

In this article, we address the corresponding recognition problem. Starting with a finite poset $P$, we determine exactly which partitions $P=A \sqcup B$ realize $P$ as the positive poset in Ladkani's construction. Our main result, Theorem \ref{generalizzazione teomio}, gives necessary and sufficient conditions for such a realization, formulated entirely in terms of the order structure of $P$. These conditions define an admissible cut (Definition \ref{def admissible cuts}): $A$ is an ideal, $B$ is a filter, and the sets $M(a)$ of minimal elements of $B$ lying above $a$ satisfy separation and compatibility conditions.

Whenever such a realization exists, the subsets required by Ladkani's construction are uniquely determined by $P$ and by the chosen partition. Moreover, such subsets coincide with the sets $M(a)$. Thus, for each partition, the construction data can be recovered from the given order and their admissibility checked directly. Testing all partitions gives a finite procedure for finding all admissible cuts of $P$ (see the discussion following Theorem \ref{generalizzazione teomio}). Each admissible cut then determines an associated poset $P^-$, obtained by changing the relations between $A$ and $B$, called mixed relations. This associated poset is universally derived equivalent to $P$ (Corollary \ref{The associated universally derived equivalent poset}). For example, applying this procedure to the five-element Tamari lattice $T_3$ produces a universally derived equivalent poset whose Hasse diagram is an orientation of the Dynkin diagram $D_5$ (Example \ref{T3/D5}). In particular, the underlying undirected Hasse graph changes from a cycle to a tree.

The transformation is reversible once the partition $A \sqcup B$ is retained: the sets $M(a)$, and hence all mixed relations of $P$, can be recovered from $P^-$ (Proposition \ref{Reversibility of the construction}). This reconstruction is dual to the original construction. We also show that every finite poset of height at most one is universally derived equivalent to its opposite (Corollary \ref{Two-level posets and opposites}).

In the singleton case, where each $M(a)$ consists of a single element, the construction is determined by an order-preserving map $A \to B$ (Corollary \ref{Reconstructing a poset from a cut}). The non-singleton case is essential, however, and appears naturally in the source-to-sink operations studied in Section \ref{An intrinsic criterion for source-clicks}. We also establish an invariance property under isomorphisms (Proposition \ref{invgen}): transporting the combinatorial data along isomorphisms produces isomorphic posets, reducing redundancy in classification problems.

We next consider local transformations of Hasse diagrams. For a minimal element $s$, we reverse all arrows incident with $s$ in the Hasse diagram and take the transitive closure, obtaining a new poset in which $s$ is maximal. This operation is called a source-click at $s$ (see Section \ref{An intrinsic criterion for source-clicks}). We give a criterion, valid for arbitrary finite posets, for this source-click to arise from an admissible cut (Proposition \ref{General source-click criterion}). For orientations of trees the criterion is automatically satisfied (Corollary \ref{Source step}). We then give a short combinatorial proof of the known fact that any orientation of a finite tree can be reached from any other by a finite sequence of source-clicks (Proposition \ref{Transitivity of clicks on a tree}). In this way, the general construction recovers the universal derived equivalence of all orientations of a tree, while also explaining it through explicit local transformations (Theorem \ref{tree-exhaustiveness}).

Finally, we relate these constructions to persistence modules, which are a central objects in topological data analysis (TDA). Since persistence modules indexed by a finite poset are diagrams of vector spaces over that poset, universal derived equivalence induces derived equivalences between the corresponding categories of pointwise finite-dimensional persistence modules (see Section \ref{Application to persistence modules and parameter extension}). We also use the stability of universal derived equivalence under products (Proposition \ref{Stability under parameter extension}) to adjoin further finite parameters to the pairs produced by admissible cuts (Corollary \ref{Parameter extension for admissible cuts}). This yields infinite families of derived equivalent parameter posets (Example \ref{example Tamari TDA}) and extends the tree-orientation results to multiparameter settings (Corollary \ref{Parameter extension for tree orientation}).

\subsection*{Organization.}
Section 2 introduces the notation and terminology used throughout the paper and recalls Ladkani's construction. Section 3 proves the main characterization theorem, shows that the associated construction is reversible, derives an application to posets of height at most one and their opposites, discusses the singleton specialization and cardinality constraints, and records an invariance property under isomorphisms. Section 4 gives the general criterion for source-clicks to be realized by admissible cuts and then applies this criterion to orientations of finite trees. Section 5 relates these results to persistence modules and studies their stability under parameter extension, with applications to the Tamari example and to orientations of finite trees.